% ============================================================================
% Chapter 6 --- Identity and subject trust
% ----------------------------------------------------------------------------
% Purpose : functional layer covering subjects, attributes, role activation
%           and federation. Belongs to the Subject and Trust plane (D5.1).
%           First layer chapter --- instantiates the per-layer template that
%           Ch. 7--14 will mirror.
% Page target: ~5 pages.
% Dependencies: Ch. 3 (sovereignty cases must remain consistent), Ch. 5
%               (standards table, replaceability conditions, three-plane
%               taxonomy).
% ----------------------------------------------------------------------------
% Decisions registered in _coherence-EN.md:
%   D6.1 minimum compliance set for an identity provider (the contract)
%   D6.2 reference instantiation = Keycloak + SCIM + eduGAIN
%   D6.3 alternatives are commercial (Entra) and hybrid (Shibboleth + Keycloak)
%   D6.4 sovereignty scores in this chapter MUST match Ch. 3 §3.4
% ============================================================================

\chapter{Identity and subject trust}
\label{ch:identity}

\begin{caveatbox}[Dependencies]
\noindent This chapter applies the architectural principles of
\trchap{ch:architectural-principles}{} to the identity layer. It
references the standards table of \S\ref{sec:ap-standards} and the four conditions of
replaceability of \S\ref{sec:ap-replaceability}. The sovereignty profiles given in
\S\ref{sec:id-sovereignty} are consistent with the worked examples of
\trchap{ch:sovereignty-grid}{} (Cases 1--3). The layer belongs to the
\emph{Subject and trust plane}.
\end{caveatbox}

The identity layer answers the first question of any access decision:
\emph{who is acting, with what attributes, and under what session
guarantees}. It is the layer on which Policy
(\trchap{ch:policy}{}), Endpoint (\trchap{ch:endpoint}{}) and
Cryptographic services (\trchap{ch:crypto}{}) all rely; a deficiency
here propagates to the entire control plane.

%-----------------------------------------------------------------------------
\section{Functional role and interface contract}
\label{sec:id-contract}

The identity layer is responsible for: authenticating subjects
(human and non-human); maintaining the canonical attributes attached
to each subject (affiliation, role, group, posture claim); issuing
short-lived assertions that downstream layers can verify; and,
where such arrangements are available and relevant to the
institution, participating in federation arrangements that allow it
to delegate authentication to peer institutions and accept their
delegated authentications in return.

\begin{contractbox}[Identity layer]
\noindent A compliant identity provider, in the sense of this
architecture, must:

\begin{itemize}[leftmargin=1.5em,nosep]
  \item issue authentication assertions over OIDC
    Core~1.0~\cite{oidc-core-1.0} \emph{or} SAML~2.0~\cite{saml-core-2.0}
    (legacy compatibility), without captive extensions in the
    assertion path;
  \item accept and emit identity provisioning over
    SCIM~2.0~\cite{rfc-7643-scim-core,rfc-7644-scim-protocol};
  \item be capable of publishing federation metadata interoperable
    with \loc{identity-federation-profile}, and do so
    where the institution participates in inter-institutional
    federation;
  \item express attributes in a documented schema, with an
    attribute-release policy reviewable by the institution;
  \item support the addition of a multi-factor authentication step
    that the institution can configure without supplier intervention;
  \item produce an authentication and assertion log in a form
    ingestible by the observability layer
    (\trchap{ch:observability}{}).
\end{itemize}
\end{contractbox}

The contract is deliberately silent on \emph{which} attribute schema
is used, on \emph{which} MFA factor is enforced, and on \emph{which}
research-and-education identity federation profile the
institution participates in. These are matters of institutional
context; the contract requires only that the answers be documented
and that the underlying mechanism be standards-based. An institution
with no national research-and-education identity federation available
satisfies the contract by
retaining the capability to publish interoperable metadata and
deferring actual participation until it joins one --- the standards
basis is required at all scales, the federation membership only
where it exists.

%-----------------------------------------------------------------------------
\section{Reference instantiation: open-source IdP + SCIM + research-and-education federation}
\label{sec:id-reference}

The reference instantiation comprises a self-hosted open-source
identity-provider deployment --- Keycloak~\cite{keycloak} is one
mature representative --- a \gls{scim} provisioning bus,
and participation in \loc{identity-federation-profile}~\cite{edugain}
through its national research-and-education federation operator.

\begin{instantiationbox}[Identity layer]
\noindent\textbf{Topology.} An open-source identity-provider realm
(here illustrated with Keycloak~\cite{keycloak}) hosts the
institutional subject base; a SAML2 metadata feed integrates with the
national research-and-education federation, which is itself a member of
the wider interfederation~\cite{edugain}; SCIM 2.0 provisioning
synchronises authoritative subject attributes into the identity
provider from the
institution's source-of-truth (typically the human resources system
for staff and the student information system for students); OIDC and
SAML2 are exposed to applications, with OIDC preferred for
modern integrations. Authentication and assertion logs are forwarded
in OpenTelemetry format to the observability layer.
\end{instantiationbox}

\paragraph{Operational considerations.} The reference instantiation
is expected to sustain the identity needs of a research-intensive
institution within a small operating team sized to the institution's
own staffing profile --- provided that the lifecycle
automation (SCIM upstream and downstream) is robust. Actual envelopes
vary with the application portfolio and the federation
participation. The principal recurring effort is the
maintenance of attribute-release policies as new applications join
the federation; this effort scales with application count, not user
count.

\paragraph{Federation profile.} The applicable
research-and-education federation's technical profile of
SAML2 is co-normative with the standard itself: the institution must
publish metadata signed by its national research-and-education
federation operator, must commit to
attribute release commitments encoded in the federation's policy
framework, and must respond to deprecation cycles (typically
multi-year) for cryptographic algorithms.

\begin{realworldbox}[CERN SSO]
\noindent CERN operates a Keycloak-based federated SSO~\cite{cern-keycloak-2020} serving its
research community and a wide set of collaborating institutions.
Aspects of the architecture have been presented in the GÉANT
community and illustrate how a research-intensive institution may
sustain a fully open identity plane without sacrificing operational
depth. The CERN experience also illustrates a recurring pattern:
the migration to Keycloak from the prior bespoke system was reported
not as a single-step replacement but as a multi-year coexistence
during which both stacks remained authoritative for distinct
application sets. Specific scale figures and migration durations
should be confirmed against current public CERN documentation before
being cited in derivative work.
\end{realworldbox}

%-----------------------------------------------------------------------------
\section{Alternative~1 --- Incumbent commercial identity suite with conditional-access policy}
\label{sec:id-alt1}

\begin{alternativebox}[Incumbent commercial identity suite]
\noindent A tenancy of \loc{incumbent-identity-suite-vendor}~\cite{microsoft-entra}, operated by the supplier in a
designated EU (or analogous regional) region, integrated through the
institution's existing tenant arrangements, with the vendor's
conditional-access policy engine
expressing the role-activation rules. Authentication is exposed to
applications via OIDC and SAML2; provisioning is exposed via SCIM
2.0; logs are forwarded to the observability layer through the
vendor's management API
or via the OpenTelemetry collectors of
\trchap{ch:observability}{}.
\end{alternativebox}

\paragraph{Where it adds value.} The operational dimension benefits
from a deep and broadly available skills market, extensive
documentation, and integration with an incumbent office-and-collaboration
suite, where one is already dominant in the institution. The cost
of activation, in operational time, is lower than for a self-hosted
solution.

\paragraph{Where it constrains the institution.} Three constraints
recur. First, the vendor's conditional-access expressions, while
powerful, are
not portable: an authorisation logic encoded in that vendor-specific
policy language
does not transfer to another identity provider without rewriting.
Second, several capabilities of operational interest (privileged-access
management, verifiable-credential issuance, advanced risk analytics) are
gated behind specific licences and may shift between licence tiers
across renewals. Third, the assertion and provisioning paths admit
optional extensions whose use would silently couple the institution
to the supplier; an architecture-aware deployment pins the
configuration to the standard interface and forgoes those extensions
in critical paths.

\paragraph{When it is the right choice.} An institution with a
mature deployment of an incumbent office-and-collaboration suite,
limited identity-engineering capacity,
and a sovereignty position that accepts a Partial score on data,
technical and institutional dimensions in exchange for a Robust score
on operational. The choice is consistent with this architecture
\emph{provided} that the OIDC, SAML2, and SCIM interfaces are pinned
to the standard and that an exit-data export procedure is documented
and rehearsed.

%-----------------------------------------------------------------------------
\section{Alternative~2 --- Hybrid (open-source federation engine + open-source local activation)}
\label{sec:id-alt2}

\begin{alternativebox}[Hybrid identity]
\noindent An institution that already operates a mature open-source
federation engine --- Shibboleth~\cite{shibboleth} is a common
representative --- preserves it for
inter-institutional authentication (facing the research-and-education
federation), and adds a
self-hosted open-source activation identity provider that
handles \emph{local activation}: role enrichment, TAC session
issuance, attribute mapping toward downstream applications, and
short-lived OIDC token emission. The two components communicate
through standard SAML2 / OIDC bridges; neither is captive to the
other.
\end{alternativebox}

\paragraph{Why this pattern works.} It separates two concerns that
have different lifetimes and different governance bodies. Federation
is long-lived, governed by inter-institutional agreements, and
benefits from the stability of a mature open-source federation engine.
Activation is short-lived,
local, and benefits from the policy expressiveness and modern
extensibility of an open-source activation identity provider. Each
component speaks the standards of
\S\ref{sec:ap-standards}; the bridge between them is the authentication context emitted
by the federation engine and consumed by the activation provider.

\paragraph{Operational implications.} The hybrid pattern carries a
modest operational premium over a
single-component instantiation, sized to the institution's staffing
profile and divided across two skill sets
(federation-engine and activation-provider operators). It is offset by
the preservation
of the existing federation investment and by the avoidance of a
high-risk single migration.

\paragraph{When it is the right choice.} An institution with a
significant existing investment in an open-source federation engine, a
research-intensive profile that already values open-source identity,
and a pace constraint that does not permit a wholesale migration in
the immediate term. The pattern is also a stable end-state, not only
a transition.

%-----------------------------------------------------------------------------
\section{Sovereignty grid applied}
\label{sec:id-sovereignty}

The three instantiations above were profiled in
\trchap{ch:sovereignty-grid}{} (Cases~1--3) and the scores are
reproduced here for the layer chapter, with brief commentary at the
indicator level.

\begin{table}[H]
\centering\small
\caption{Per-dimension sovereignty profile of the three identity
instantiations. Scores are consistent with
\trchap{ch:sovereignty-grid}{} \S\ref{sec:sg-examples}.}
\label{tab:id-sovereignty}
\begin{tabularx}{\textwidth}{%
  >{\hsize=0.80\hsize\linewidth=\hsize\bfseries\raggedright\arraybackslash}X%
  >{\hsize=0.80\hsize\linewidth=\hsize\centering\arraybackslash}X%
  >{\hsize=0.80\hsize\linewidth=\hsize\centering\arraybackslash}X%
  >{\hsize=0.80\hsize\linewidth=\hsize\centering\arraybackslash}X%
  >{\hsize=1.80\hsize\linewidth=\hsize\raggedright\arraybackslash}X%
}
\toprule
Dimension & Pure OSS & Commercial suite (default) & Hybrid & Driving indicators \\
\midrule
Data          & 3 Robust      & 1 Partial      & 3 Robust      & Storage location, key custody, disclosure regime \\
Technical     & 3 Robust      & 1 Partial      & 3 Robust      & Source visibility, standard interfaces, regression auditability \\
Financial     & 2 Established & 1 Partial      & 2 Established & Cost predictability, supplier substitutability \\
Operational   & 2 Established & 3 Robust       & 2 Established & Local skills availability, labour-market depth \\
Institutional & 3 Robust      & 1 Partial      & 3 Robust      & Roadmap influence, integration freedom \\
\bottomrule
\end{tabularx}
\end{table}

\noindent The principal sovereignty trade-off in the identity layer
is between the operational depth of the commercial alternative and
the technical / institutional / data sovereignty of the open
alternatives. The hybrid pattern recovers most of the open-side
profile while modestly increasing the operational footprint over a
single-component instantiation, by an increment sized to the
institution's own staffing profile.

%-----------------------------------------------------------------------------
\section{Regulatory anchoring}
\label{sec:id-regulatory}

\noindent The identity layer is the primary technical anchor of
the access-control security obligations of the applicable
data-protection regime (e.g.\ \gls{gdpr}~Art.~32(1)(b)) and of
several measures of the applicable cybersecurity baseline regime
(e.g.\ NIS2~Art.~21(2)): (i) human-resources security and
access control, (j) multi-factor authentication and secured
communications. Joiner / mover / leaver via SCIM operationalises
27001~A.5.16, A.6.5 and A.8.2; the OIDC / SAML2 federation pattern
operationalises eIDAS~2 trust services and notified-eID
interoperability; the MFA / FIDO2 baseline operationalises
27001~A.8.5. The layer contributes to the NIST CSF \textsf{PROTECT}
function (PR.AA --- identity management, authentication and access
control). Detailed mapping in \trchap{ch:regulatory-mapping}{}
\S\ref{sec:rm-gdpr}, \S\ref{sec:rm-nis2}, \S\ref{sec:rm-eidas2},
\S\ref{sec:rm-iso}, \S\ref{sec:rm-nist-csf}.

%-----------------------------------------------------------------------------
\section{Common pitfalls}
\label{sec:id-pitfalls}

\begin{pitfallbox}[Captive attribute extensions]
\noindent Encoding authorisation logic in supplier-specific attribute
expressions (vendor conditional-access expressions, vendor-defined claims)
turns the identity provider into a captive policy engine. The
mitigation is to keep authorisation in the policy layer
(\trchap{ch:policy}{}) and to use the identity layer only for
authentication and attribute carriage, against the standard schema.
\end{pitfallbox}

\begin{pitfallbox}[Broken metadata refresh]
\noindent Federation metadata expires, is rotated, and is signed by
keys that themselves rotate. Institutions whose metadata refresh
pipeline is undocumented or manually triggered lose federation
participation silently --- typically discovered during a partner
authentication outage. The mitigation is to automate metadata
refresh, monitor freshness through the observability layer, and
include a metadata-currency check in the conformance suite of
\trchap{ch:conformance}{}.
\end{pitfallbox}

\begin{pitfallbox}[Opaque multi-factor risk decisions]
\noindent Some commercial identity providers couple MFA to a
proprietary risk score whose inputs and weights are not disclosed.
The institution then cannot explain to its own users why a session
was challenged, nor reproduce the decision for audit. The mitigation
is to either use a transparent MFA factor (TOTP, FIDO2 with documented
attestation handling) or, if a risk-based factor is used, to ensure
that the decision is logged with sufficient detail for explanation.
\end{pitfallbox}

\begin{pitfallbox}[Attribute-release drift]
\noindent The set of attributes released to each application tends to
grow over time, and to lose alignment with what the application
actually requires. This is both a data-minimisation issue under the
applicable data-protection regime and a sovereignty issue, since unnecessary releases
multiply the dependencies that an exit must address. The mitigation
is an annual review of the attribute-release policy by application,
recorded in the institution's data-protection register.
\end{pitfallbox}
